\section{The ordinary realization and the bridge to iPRA}
\label{sec:new-bridge}

Here ``iPRA'' means full intuitionistic first-order logic with actual
numerical equality, a function symbol for every fixed finite primitive
recursive description, the defining equations, successor separation
and injectivity, and induction for quantifier-free formulas with
arbitrary retained assumptions.  It does not mean just its equational
fragment, and no excluded-middle principle for quantified formulas
is added.

The existing scalar realization $\Teq$ uses ordinary intuitionistic
logic on $0,S,I$, the two iterator equations, the coverage axiom
\[
           \forall x\,(x=0\lor\exists y\,x=Sy),
\]
the falsity bridge $(0=S0)\to\bot$, and contextual induction for
literal equations.  A fixed update term is elaborated into a fixed
function description.  There is no numerical variable ranging over
programs and no assumed universal primitive recursive evaluator.

\subsection{The composed interpretations}
Let $j:\Mcal\to\Ccal$ be inclusion of the canonical numerical and
formula syntax, and let $T:\Ccal\to\Mcal$ be the padding translation
proved above.  Let $R$ and $U$ be the inherited interpretations from
$\Ccal$ to iPRA and back.  The new pair is
\[
                  Rj:\Mcal\longrightarrow\ipra,
       \qquad TU:\ipra\longrightarrow\Mcal.
\]
The preceding proofs give $TjB\dashv\vdash B$ and
$jTA\dashv\vdash A$.  The inherited formula composites give
$UR A\dashv\vdash A$ and $RU\varphi\leftrightarrow\varphi$.
Applying the faithful translations and composing these equivalences
yields both new formula composites.  This records the actual
bi-interpretation argument; mutual interpretability alone would not
have been enough.

The numerical maps between $\Mcal$ and $\Ccal$ are identities.
Between $\Teq$ and iPRA they replace each fixed function description
by a finite scalar term or vice versa.  Their internal round trips
are proved by uniqueness of primitive recursion.  The same whole
numerical domains and identity graphs therefore serve for the new
composites.  Native $0=S0$ translates to ordinary falsity in both
directions, so finite proof translations give equiconsistency.

\subsection{Why the scalar iterator has the required arithmetic}
The numerical argument is inherited, not a new assumption hidden in
the proposed proof checker.  The complete written sources are
reproduced in the appendices.  Their dependency order is as follows.
\begin{enumerate}
\item Literal-equation induction and numerical cases construct a
scalar predecessor $P$ with $P0=0$ and $P(Sn)=n$.  This is the part
that cannot be replaced by a casual assertion that an iterator can
simulate arbitrary recursion: the update has only one scalar state
and has no recursion-counter binder.
\item Predecessor gives successor injectivity.  Selector computations
and separation give constructor disjointness.  Truncated subtraction
and an equality test are then derived and proved to reflect actual
equality throughout the domain.
\item Characteristic terms for quantifier-free formulas yield full
quantifier-free induction, including its intuitionistic propositional
connectives.  The equality test is proved internally before it is used.
\item Pairing and its two projections are constructed using bounded
scalar iteration.  General primitive recursion is compiled into an
iterator storing both its counter and its previous value.
\end{enumerate}
For example, with a proved pairing $\pi$ and projections $L,H$, a
target recursion
\[
 f(0,\vec p)=g(\vec p),\qquad
 f(Sn,\vec p)=h(n,f(n,\vec p),\vec p)
\]
is represented by
\[
 V(n,\vec p)=I\bigl(n,\pi(0,g^*(\vec p));
   u.\pi(SL(u),h^*(L(u),H(u),\vec p))\bigr),
 \qquad f^*(n,\vec p)=H(V(n,\vec p)).
\]
Equation induction proves $L(V(n,\vec p))=n$, with all parameters
fixed.  The two projection identities then prove both defining
equations for $f^*$.  Uniqueness of the two recursive constructions
proves the internal function composites, including at nonstandard
elements of models; standard-model extensional agreement is not being
substituted for that argument.

\subsection{Exact inherited formal statements}
The pinned file \texttt{NativeBiInterpretation.lean} contains
\texttt{native\_iff\_pra}, \texttt{pra\_iff\_native},
\texttt{native\_roundTrip}, \texttt{pra\_native\_roundTrip}, and
\texttt{native\_inconsistent\_iff}.  The scalar bridge is
\texttt{biInterp : BiInterp TeqAx PRAAx} in
\texttt{BiInterpretation.lean}.  The reports' rule presentation is
connected to that native kernel by \texttt{iprf\_iff\_nprf}.
These are the inherited results used here, rather than a claim that
the new mirror restriction has already been inserted into Lean.

A Lean continuation would introduce the mirror-guard predicate,
the root-determined proof relation, and the two syntax translations,
then formalize safe right expansion, shared-left implication, the
substitution-free induction construction, and the formula composites.
It should not start by restricting the existing transport constructor
and invoking its structural substitution theorem: a source variable
can become the forbidden literal $0$ under such a substitution.
After full transport has been recovered, unrestricted substitution
is admissible through the established equivalence.
